\documentclass[11pt,a4paper]{article}
\usepackage[top=2cm,bottom=2cm,left=2.2cm,right=2.2cm]{geometry}
\usepackage[T1]{fontenc}
\usepackage{lmodern}
\usepackage{microtype}
\setlength{\parindent}{0pt}
\setlength{\parskip}{6pt plus 1pt minus 1pt}
\pagenumbering{gobble}

\begin{document}

\begin{flushright}
Hern\'an D\'iaz Rodr\'iguez\\
University of Oviedo\\
Department of Computer Science\\[2pt]
June 6, 2026
\end{flushright}

\bigskip

\noindent
The Editors\\
\emph{Computers \& Operations Research}\\
Elsevier

\bigskip

\noindent Dear Editor,

I am pleased to submit my manuscript, ``\textbf{Neighbourhood Structures and
Ranking Operators for the Interval Job Shop Scheduling Problem},'' for
consideration as a regular research article in \emph{Computers \& Operations
Research}. The paper addresses the Interval JSP, where processing times are
bounded but uncertain, and contributes on both theoretical and computational
fronts.

On the theoretical side, I formalise the notion of \emph{extreme critical
paths} --- arcs critical in either the best- or worst-case scenario graph
--- and define a neighbourhood $H(\sigma)$ on this basis. I prove that
$H(\sigma)$ satisfies feasibility, connectivity and the
no-loss-of-improving-neighbours property under any admissible interval
ranking operator. Five neighbourhood variants ($N_1$, $N_2$, $N_3$,
$N_{\rm ext}$ and a new variant $N_8$) are formulated within this framework;
the adaptations of the Nowicki--Smutnicki extended neighbourhood and of
recent crisp-JSP reinsertion moves to the interval setting are, to the best
of my knowledge, novel. On the empirical side, the paper reports a
$12{,}300$-run study on 82 benchmark instances of up to 1000 operations,
organised in three phases: a $5{\times}4$ comparison of four interval
ranking operators (EV, LEX1, LEX2, YX), a head-to-head comparison of the
five neighbourhoods after per-neighbourhood irace tuning of a tabu search
metaheuristic, and a comparison against three published IJSP solvers. The
irace-tuned tabu search with $N_2$ reduces mean relative error by
$\approx$57\% over the previous best published IJSP solver on every one of
the 82 instances.

I believe the paper aligns well with the scope of \emph{Computers \& Operations
Research}: it combines a constructive theoretical framework with an extensive
numerical study and explicit comparisons against the prior state of the art.
In line with the journal's emphasis on reproducibility, all raw run logs,
aggregated CSVs, irace tuning records, instance files and analysis scripts are
openly available on Zenodo (\texttt{10.5281/zenodo.20562888}), with per-instance
results on the 70 Taillard instances in a separate supplementary document
(\texttt{10.5281/zenodo.20562857}). I confirm that this manuscript has not been
published previously, is not under consideration by any other journal, and
that all listed contributors have approved the submission. I declare no
competing financial interests or personal relationships that could have
appeared to influence the work reported. Thank you for considering this
submission.

\bigskip

\noindent Sincerely,

\medskip

\noindent
Hern\'an D\'iaz Rodr\'iguez\\
University of Oviedo\\
Department of Computer Science\\
\texttt{diazhernan@uniovi.es}

\end{document}
